We take an early state vector from a Sentinel-1 precise-orbit (EOF) file as
the initial condition, propagate one hour, and RMSE-compare against the EOF
truth (RTN-decomposed) via \code{compute\_propagation\_result}. We sweep the
integrator (\code{RK45} adaptive vs \code{YOSHIDA8} symplectic) and the ITRF
realization.

\textbf{ITRF rule.} ESA switched Sentinel POD products from ITRF2014 to
ITRF2020 on 4 March 2025. The 2016 data here is \textbf{ITRF2014}; using
ITRF2020 is the ``wrong frame'' contrast.

\examplenote{\code{examples/cpp/ex03\_sentinel\_propagation.cpp}}
\exfile{examples/cpp/ex03_sentinel_propagation.cpp}{C++}
\resultfile{results/ex03_sentinel_rmse.json}{1-hour RMSE vs EOF truth.}

RK45/ITRF2014 lands at ${\sim}0.18\,$m --- comfortably inside the
documented sub-meter Sentinel baseline (the 6-satellite January-2026
campaign gives $0.54\,$m with these bare defaults and $0.24\,$m with the
recommended configuration; this early-2016 arc is a friendly one). Two
library fixes moved it: the 2026-07 epoch-label/interpolation fix (from
$0.41\,$m) and the 2026-11 exact \code{Epoch}, which removed a
$13\,\mu$s short landing of every one-hour target (along-track
$0.10 \to 0.02\,$m). YOSHIDA8 is competitive on this arc.
The ITRF2014-vs-ITRF2020 difference is a few mm: real, but second-order
next to frame and EOP correctness --- use the right realization anyway.
